
\documentclass[11pt,a4paper]{article}

\usepackage[margin=1in]{geometry}
\usepackage[utf8]{inputenc}
\usepackage[T5]{fontenc}
\usepackage[vietnamese,english]{babel}
\usepackage{lmodern}
\usepackage{newunicodechar}
\newunicodechar{✓}{\ensuremath{\checkmark}}
\newunicodechar{✗}{\ensuremath{\times}}
\usepackage{amsmath,amssymb}
\usepackage{booktabs}
\usepackage{tabularx}
\usepackage{array}
\usepackage{caption}
\usepackage{hyperref}
\usepackage{enumitem}
\usepackage{microtype}
\usepackage{graphicx}
\usepackage{float}
\usepackage[ruled,vlined,linesnumbered]{algorithm2e}
\SetKwInput{KwIn}{Đầu vào}
\SetKwInput{KwOut}{Đầu ra}
\SetKwProg{Fn}{Hàm}{}{}
\renewcommand{\algorithmcfname}{Thuật toán}

\setlength{\parindent}{0pt}
\setlength{\parskip}{0.65\baselineskip}
\hypersetup{colorlinks=true,linkcolor=blue,citecolor=blue,urlcolor=blue}
\captionsetup{font=small,labelfont=bf}

\title{Adaptive and Intelligent Data Block Scheduling for Fault-Tolerant Edge Data Distribution}
\author{Nguyen Minh Tri, Nguyen Ngoc Thong\\Trường Đại học Công nghệ thông tin}
\date{}

\begin{document}
\maketitle

\begin{abstract}
Trong bối cảnh điện toán biên (Edge Computing) ngày càng phát triển, việc phân phối khối lượng dữ liệu khổng lồ từ trung tâm đám mây (Cloud) xuống các trạm biên (Edge Nodes) đang trở thành một thách thức lớn, đặc biệt khi yêu cầu về độ trễ và độ tin cậy được đặt lên hàng đầu. Đồ án này tập trung nghiên cứu, hiện thực hóa và đánh giá chuyên sâu một giải pháp phân phối dữ liệu chịu lỗi mạng biên mang tên EdgeHydra, dựa trên nền tảng kỹ thuật mã hóa xóa (Erasure Coding), nhằm khắc phục triệt để những điểm nghẽn chí mạng của mô hình truyền dữ liệu song song truyền thống (baseline EdgeDis). 

Về mặt triển khai, nhóm tác giả đã thiết kế và xây dựng từ đầu một hệ thống mô phỏng mạng phân tán hoàn chỉnh sử dụng ngôn ngữ Python với kiến trúc lập trình bất đồng bộ (asyncio), kết hợp cùng thư viện \textit{zfec} để áp dụng quy tắc mã hóa xóa với tỷ lệ (k=3, m=1). Toàn bộ hệ sinh thái mạng bao gồm 1 máy chủ Cloud và 4 máy chủ Edge được tự động hóa và đóng gói thông qua nền tảng Docker Compose, đảm bảo môi trường thực thi độc lập và sát với thực tế. Điểm nổi bật của đồ án là việc tích hợp một bảng điều khiển Web Dashboard (FastAPI), cho phép giám sát trực quan cấu trúc liên kết mạng (Topology) và tiến trình truyền tải theo thời gian thực. Để thử nghiệm giới hạn chịu đựng của hệ thống, nhóm đã phát triển một Bộ giả lập sự cố (Fault Simulator) linh hoạt, có khả năng tiêm (inject) các lỗi mạng đặc thù như nghẽn mạng cục bộ (Slow Servers) hay sập node đột ngột (Failed Servers). 

Hơn thế nữa, nghiên cứu này không chỉ dừng lại ở việc tái hiện thuật toán gốc mà còn đóng góp một cải tiến độc quyền mang tính đột phá: Thuật toán lập lịch truyền thích ứng (Adaptive Scheduling) kết hợp kiểm soát độ phủ (Controlled Fanout). Bằng cách liên tục đo lường độ trễ khứ hồi (EWMA RTT) và tỷ lệ thành công của các kết nối ngang hàng, hệ thống tự động điều hướng dữ liệu tối ưu, qua đó loại bỏ hoàn toàn hiện tượng "bão tin nhắn" (Broadcast Storm) vốn rất phổ biến trong các kiến trúc chia sẻ phi tập trung (Leaderless).

Kết quả đo đạc thực nghiệm trên nhiều kịch bản đa dạng với kích thước tệp dữ liệu lên tới 144 MiB đã chứng minh tính ưu việt vượt trội của giải pháp đề xuất. Trong các tình huống mạng xảy ra sự cố nghiêm trọng, EdgeHydra thể hiện sự ổn định tuyệt đối, duy trì năng lực phân phối dữ liệu với độ trễ thấp hơn trung bình 10 lần (tốc độ cải thiện lên tới 1000\%) so với mô hình EdgeDis. Sự đánh đổi (trade-off) về mức độ gia tăng chi phí băng thông chuẩn hóa (Transmission Cost) để đổi lấy khả năng dung lỗi này được kiểm soát ở mức rất thấp, chỉ chênh lệch khoảng 7\% đến 16\%. Thông qua kết quả thực nghiệm toàn diện, đồ án đã khẳng định tiềm năng ứng dụng mạnh mẽ của EdgeHydra cho các hạ tầng Edge Caching đòi hỏi độ trễ siêu thấp và khả năng mở rộng (Scalability) cao trong tương lai.
\end{abstract}


\textbf{Keywords:} Edge Computing; Edge Data Distribution; Erasure Coding; Fault Tolerance; EdgeHydra; EdgeDis.


\section{Giới thiệu}

Sự bùng nổ của mạng di động thế hệ mới (5G/6G), sự gia tăng theo cấp số nhân của các thiết bị Internet vạn vật (IoT), và nhu cầu ngày càng cao của các ứng dụng trí tuệ nhân tạo (AI) thời gian thực như xe tự lái hay thực tế ảo (VR/AR) đã đặt ra những yêu cầu vô cùng khắt khe về băng thông và độ trễ phản hồi ở mức mili-giây. Để giải quyết bài toán này, kiến trúc tính toán biên (Edge Computing) kết hợp với kỹ thuật lưu trữ tạm thời tại biên (Edge Caching) đã ra đời như một mô hình mang tính cách mạng. Bằng cách kéo dữ liệu phổ biến từ các trung tâm dữ liệu đám mây (Cloud) khổng lồ về lưu trữ tại các trạm biên (Edge Servers) nằm sát cạnh người dùng cuối, giải pháp này không chỉ giúp giảm thiểu đáng kể thời gian truy xuất dịch vụ mà còn giảm tải áp lực nghẽn mạng cho hạ tầng đường trục (backhaul network).

Tuy nhiên, việc phân phối một lượng dữ liệu lớn (như các mô hình Deep Learning hoặc các tệp video chất lượng cao) từ Cloud xuống mạng lưới các node biên lại đang đối mặt với hai thách thức lớn về chi phí kinh tế (Cost) và độ tin cậy của mạng (Reliability). Thứ nhất, các nhà cung cấp dịch vụ đám mây (như AWS, Google Cloud) thường tính phí băng thông truyền tải dữ liệu từ Cloud xuống Edge (egress traffic) đắt hơn gấp 5 đến 9 lần so với chi phí trao đổi dữ liệu nội bộ ngang hàng (P2P LAN) giữa các Edge node với nhau. Thứ hai, hạ tầng mạng biên thường bị giới hạn về năng lực tính toán và chất lượng kết nối, khiến các trạm biên rất dễ gặp phải các sự cố bất khả kháng như nghẽn mạng cục bộ (Slow Servers), rớt gói tin, hoặc thậm chí là sập nguồn hoàn toàn (Node Crash). Các phương pháp truyền thống tỏ ra kém hiệu quả trong môi trường này: phương pháp DataSync (gửi song song toàn vẹn từ Cloud xuống mọi node) gây hao phí băng thông đám mây cực lớn, trong khi thuật toán Gossip (lan truyền ngẫu nhiên) lại có độ trễ hội tụ rất cao và khó kiểm soát.

Gần đây, giao thức EdgeDis đã được giới học thuật đề xuất như một giải pháp tối ưu chi phí bằng cách chia nhỏ tệp tin gốc thành các phân đoạn tĩnh không mã hóa, gửi một phần từ Cloud và để các trạm biên tự trao đổi chéo với nhau. Dẫu vậy, nhược điểm chí mạng của EdgeDis là tính phụ thuộc tuyệt đối vào từng phân đoạn dữ liệu riêng lẻ và sự tồn tại của một trạm biên đóng vai trò điều phối trung tâm (Coordinator). Cơ chế này vô tình tạo ra một điểm lỗi đơn (Single Point of Failure). Chỉ cần một trạm biên đóng vai trò Entry Server bị chậm hoặc sập, toàn bộ tiến trình phân phối sẽ bị ngưng trệ, buộc Coordinator phải chờ đợi hết thời gian (timeout) rồi mới yêu cầu tải lại khối dữ liệu đó từ Cloud. Hậu quả là thời gian hoàn thành bị kéo dài nghiêm trọng và lãng phí băng thông không đáng có.

Để giải quyết triệt để những điểm nghẽn về cả mặt kiến trúc lẫn thuật toán nêu trên, đồ án này tập trung nghiên cứu, hiện thực hóa và đánh giá so sánh chuyên sâu một giải pháp phân phối dữ liệu mạng biên chịu lỗi mang tên EdgeHydra, hoạt động trên nền tảng toán học của Mã hóa xóa (Erasure Coding). Bằng việc áp dụng mã hóa xóa với cấu hình phân mảnh tối ưu (k=3, m=1), EdgeHydra tạo ra các khối dữ liệu dư thừa, cho phép mỗi trạm biên chỉ cần thu thập đủ 3 khối bất kỳ (từ bất kỳ nguồn nào) là đã có thể tự giải mã để khôi phục tệp tin gốc mà không cần phải chờ đợi các node đang bị lỗi. Về mặt kỹ thuật, nhóm tác giả đã xây dựng thành công một hệ thống mạng phân tán mô phỏng bằng ngôn ngữ Python (kiến trúc asyncio) kết hợp thư viện toán học \textit{zfec}. Mạng lưới 5 node (1 Cloud, 4 Edge) được container hóa độc lập bằng Docker Compose, đi kèm một Web Dashboard (FastAPI) đóng vai trò giám sát và trực quan hóa tiến trình truyền tải theo thời gian thực.

Tóm lại, những đóng góp cốt lõi của đồ án này bao gồm:

\begin{itemize}
    \item \textbf{Hiện thực hóa hoàn chỉnh và đóng gói hệ thống phân tán:} Lập trình giả lập thành công toàn bộ luồng hoạt động song song của giao thức chịu lỗi EdgeHydra (bao gồm 4 giai đoạn truyền tin phức tạp) và giao thức baseline EdgeDis trên một môi trường mạng phân tán container hóa thực tế, không dừng lại ở mức mô phỏng toán học lý thuyết.
    
    \item \textbf{Tích hợp cơ chế Lập lịch Thích ứng (Adaptive Scheduling) và Phi tập trung (Leaderless):} Đồ án cải tiến mạnh mẽ thuật toán gốc bằng cách triển khai cơ chế kiểm soát độ phủ (Controlled Fanout) và lập lịch ưu tiên dựa trên bộ lọc trung bình trượt lũy thừa độ trễ khứ hồi (EWMA RTT) kết hợp cùng tỷ lệ thành công mạng. Kết hợp cùng cơ chế bù đắp khối phi tập trung thông qua UDP Heartbeat chu kỳ 100ms, cải tiến này giúp hệ thống loại bỏ triệt để vấn đề "Bão broadcast" (Broadcast Storm), tối ưu hóa lưu lượng mạng ngang hàng.
    
    \item \textbf{Đánh giá thực nghiệm định lượng độ tin cậy và khả năng mở rộng:} Thiết lập Bộ giả lập sự cố (Fault Simulator) linh hoạt để đo đạc hiệu năng với dữ liệu lớn (lên tới 144 MiB). Kết quả thu được chứng minh EdgeHydra mang lại tốc độ phân phối đột phá, \textbf{nhanh hơn trung bình 10 lần (cải thiện độ trễ lên tới 1000\%)} so với EdgeDis trong điều kiện xảy ra sập node đột ngột, với sự đánh đổi chi phí băng thông chuẩn hóa nội bộ (Cost Units) cực kỳ nhỏ (chỉ chênh lệch từ 7\% đến 16\%).
\end{itemize}


\section{Các nghiên cứu liên quan}

\subsection{Các phương pháp phân phối dữ liệu tại biên}

Trong những năm gần đây, nhiều giải pháp phân phối dữ liệu tại biên (Edge Data Distribution - EDD) đã được nghiên cứu nhằm tối ưu hóa chi phí băng thông và tốc độ phổ biến nội dung.

Phương pháp DataSync áp dụng chiến lược đẩy song song, tức là Cloud Server trực tiếp gửi toàn bộ nội dung tới từng node biên một cách độc lập và đồng thời. Phương pháp này đảm bảo tất cả các node nhận được nội dung đầy đủ trong thời gian ngắn, tuy nhiên chi phí băng thông Cloud-to-Edge rất cao do cùng một dữ liệu phải được truyền lặp lại nhiều lần từ Cloud.

Phương pháp Gossip tiếp cận theo hướng ngược lại: Cloud chỉ gửi một bản sao dữ liệu tới một node biên ban đầu, sau đó dữ liệu được lan truyền ngẫu nhiên theo kiểu "lây lan" giữa các node ngang hàng. Phương pháp này tiết kiệm tối đa chi phí Cloud-to-Edge nhưng dẫn đến độ trễ phân phối rất lớn và không thể kiểm soát, không phù hợp với ứng dụng yêu cầu thời gian thực.

Để khắc phục nhược điểm của hai phương pháp trên, giao thức EdgeDis được đề xuất kết hợp ưu điểm của cả hai: Cloud chia nhỏ tệp tin thành n phần bằng nhau và giao mỗi phần cho một node biên đóng vai trò Entry Server, sau đó các node tự động trao đổi chéo dữ liệu cho nhau qua mạng LAN nội bộ tiết kiệm hơn. Tuy nhiên, hạn chế chí mạng của EdgeDis là sự phụ thuộc tuyệt đối vào tất cả n phân đoạn: chỉ cần một Entry Server bị trễ mạng hoặc sập nguồn, toàn bộ quá trình phân phối tại tất cả node bị đình trệ vì thiếu đúng phân đoạn đó, buộc hệ thống phải yêu cầu tải lại từ Cloud.

\subsection{Tổng hợp và định vị đóng góp}

Từ các phân tích trên, có thể thấy rõ khoảng trống nghiên cứu chưa được giải quyết: chưa có giao thức phân phối dữ liệu tại biên nào đồng thời đáp ứng được ba yêu cầu là tiết kiệm chi phí băng thông Cloud, chịu được lỗi node biên một cách hoàn toàn tự động, và không cần điểm điều phối trung tâm (Leaderless). Bảng 1 dưới đây tổng hợp sự so sánh giữa các giải pháp hiện có và EdgeHydra.

\begin{table}[htbp]

\centering

\caption{So sánh các giải pháp phân phối dữ liệu tại biên}

\small

\begin{tabularx}{\textwidth}{>{\raggedright\arraybackslash}X>{\raggedright\arraybackslash}X>{\raggedright\arraybackslash}X>{\raggedright\arraybackslash}X>{\raggedright\arraybackslash}X}

\toprule

\textbf{Giải pháp} & \textbf{Tiết kiệm} & \textbf{Chịu lỗi node} & \textbf{Phi tập trung} & \textbf{Mã hóa xóa} \\

\midrule

DataSync & $\times$ & \checkmark & \checkmark & $\times$ \\

Gossip & \checkmark & \checkmark & \checkmark & $\times$ \\

EdgeDis & \checkmark & $\times$ & $\times$ & $\times$ \\

EdgeHydra & \checkmark & \checkmark & \checkmark & \checkmark \\

\bottomrule

\end{tabularx}

\end{table}

Như thể hiện trong Bảng 1, EdgeHydra là giải pháp duy nhất đáp ứng đồng thời cả bốn tiêu chí đánh giá, khẳng định tính mới và giá trị thực tiễn của đồ án.

\section{Phân phối dữ liệu EdgeHydra}

\subsection{Cơ sở lý thuyết: Mã hóa xóa (Erasure Coding)}

Mã hóa xóa (Erasure Coding - EC) là kỹ thuật mã hóa dữ liệu cho phép hệ thống chịu lỗi mà không cần lưu trữ dư thừa nhiều như cơ chế nhân bản (Replication) truyền thống. Với tham số cấu hình (k, m), hệ thống chia dữ liệu gốc thành k khối dữ liệu (data shards) và sinh thêm (m - k) khối kiểm tra (parity shards), tạo thành tổng số m khối được phân tán lưu trữ. Bất kỳ tập hợp nào gồm k khối trong số m khối đều đủ để khôi phục hoàn chỉnh dữ liệu gốc, ngay cả khi (m - k) khối bị mất hoặc hỏng.

Kỹ thuật này đã được ứng dụng rộng rãi và thành công trong các hệ thống lưu trữ đám mây quy mô lớn như HDFS (Hadoop Distributed File System) và Windows Azure Storage nhằm tăng độ tin cậy dữ liệu trong khi tiết kiệm không gian lưu trữ so với nhân bản 3 lần (3-way replication). Tuy nhiên, trong tất cả các ứng dụng kể trên, mã hóa xóa chỉ được sử dụng cho mục đích lưu trữ dữ liệu tĩnh, chưa được khai thác để tối ưu hóa quá trình truyền tải dữ liệu phân tán động trong thời gian thực.

Đóng góp cốt lõi của EdgeHydra là bước đột phá trong việc lần đầu tiên áp dụng mã hóa xóa trực tiếp vào giao thức truyền tải dữ liệu tại biên, biến nó thành công cụ chủ động giúp hệ thống bỏ qua các node bị lỗi mà không cần bất kỳ sự can thiệp nào từ Cloud.

\subsection{Kiến trúc tổng quan}

Hệ thống EdgeHydra được thiết kế để phân phối các tệp tin dữ liệu dung lượng lớn từ máy chủ trung tâm (Cloud) đến mạng lưới các trạm biên (Edge Servers). Mô hình thực nghiệm giả lập của đồ án thiết lập mạng lưới gồm một máy chủ đám mây trung tâm (Cloud Node) và n = 4 máy chủ biên độc lập (e\ensuremath{_0}, e\ensuremath{_1}, e\ensuremath{_2}, e\ensuremath{_3}).. Các máy chủ biên liên kết với nhau theo cấu trúc hình lưới (mesh topology) thông qua kết nối mạng nội bộ LAN ảo (được thiết lập bằng cấu hình Docker bridge network). Mỗi Edge Node chạy một tiến trình máy chủ bất đồng bộ (EdgeServer) quản lý kho lưu trữ khối dữ liệu cục bộ, thực hiện đo đạc hiệu năng kết nối mạng của các peer lân cận và tự động trao đổi các bản tin kiểm soát cũng như dữ liệu.

\begin{figure}[H]
    \centering
    % Nhớ thay thế "images/topology.png" bằng đường dẫn thật tới ảnh bạn đã chụp
    \includegraphics[width=0.8\textwidth]{images/topology.png}
    \caption{Kiến trúc mạng tổng thể của hệ thống EdgeHydra}
    \label{fig:topology}
\end{figure}

\subsection{Giao thức truyền dữ liệu 4 giai đoạn}

Trái tim của hệ thống EdgeHydra là giao thức truyền dữ liệu chịu lỗi gồm 4 giai đoạn vận hành bất đồng bộ và phi tập trung (như minh họa tại Hình \ref{fig:4stages}):

\begin{figure}[H]
    \centering
    % Nhớ thay thế "images/4stages.png" bằng đường dẫn thật tới ảnh bạn đã chụp
    \includegraphics[width=\textwidth]{images/4stages.png}
    \caption{Sơ đồ tuần tự 4 giai đoạn phân phối dữ liệu}
    \label{fig:4stages}
\end{figure}

\textbf{Giai đoạn 1 (Cloud Distribution):} Khi có yêu cầu phân phối một tệp tin, máy chủ Cloud sẽ đọc dữ liệu và áp dụng thuật toán mã hóa xóa (Erasure Coding) sử dụng thư viện zfec với cấu hình k = 3, m = 4. File gốc được phân rã thành k = 3 khối dữ liệu gốc và tính toán thêm m - k = 1 khối kiểm tra (parity), tạo ra tổng cộng m = 4 khối dữ liệu đã mã hóa (b\ensuremath{_0}, b\ensuremath{_1}, b\ensuremath{_2}, b\ensuremath{_3}). Cloud thực hiện gửi song song từng khối mã hóa b\_i tới trạm biên tương ứng e\_i (đóng vai trò là Entry Server của khối đó). Quá trình truyền tải sử dụng khung tin TCP tự định nghĩa (Custom TCP Frame) có cấu trúc: [Độ dài Header (4 bytes)] [Header Metadata (JSON)] [Độ dài Payload (4 bytes)] [Khối dữ liệu (Bytes)]. Thông điệp truyền từ Cloud có kiểu bản tin là mesbd với nội dung JSON: \{"type": "mesbd", "block\_id": i, "file\_name": "video.mp4"\}.

\begin{algorithm}[H]
\caption{Phân phối dữ liệu từ Cloud (Stage 1)}
\KwIn{Tệp $F$, Tham số mã hóa $k, m$, Tập các trạm biên $E$}
$Blocks \leftarrow \text{ErasureEncode}(F, k, m)$\;
\For{mỗi $b_i \in Blocks$}{
    $e_i \leftarrow E[i \bmod |E|]$\;
    $mesbd(b_i) \leftarrow \text{CreateMessage(type="mesbd", data=}b_i)$\;
    Gửi $mesbd(b_i)$ tới trạm $e_i$\;
}
Chờ tín hiệu hoàn thành $mesdr$\;
\end{algorithm}

\textbf{Giai đoạn 2 (Edge Forwarding):} Ngay khi một trạm biên nhận được khối dữ liệu từ Cloud thông qua thông điệp mesbd, trạm biên đó lập tức lưu khối dữ liệu vào phân vùng bộ nhớ đệm (cache) của mình. Sau đó, node này lập tức chuyển tiếp (forward) khối dữ liệu này tới các trạm biên khác trong mạng lưới bằng thông điệp mesbt thông qua các kết nối TCP đã thiết lập sẵn. Để tối ưu hóa hiệu năng và tránh gây nghẽn băng thông do gửi tràn lan, EdgeHydra tích hợp giải thuật Lập lịch truyền thích ứng (Adaptive Forwarding). Node biên sử dụng mô-đun giám sát trạng thái mạng (NetworkStateMonitor) để xếp hạng các peer lân cận dựa trên hiệu năng đường truyền và chỉ ưu tiên chuyển tiếp khối dữ liệu tới các node có kết nối tốt nhất.

\begin{algorithm}[H]
\caption{Chuyển tiếp dữ liệu thích ứng tại Trạm biên (Stage 2)}
\KwIn{Thông điệp nhận được $msg(b_i)$, Sổ cái $L$}
\If{$b_i \notin L$}{
    Lưu $b_i$ vào cache cục bộ\;
    Thêm $b_i$ vào sổ cái $L$\;
    $Peers \leftarrow \text{GetNeighborPeers}()$\;
    \ForEach{$p \in Peers$}{
        $Score(p) \leftarrow \text{EWMA\_RTT}(p) + \alpha \times (1 - \text{SuccessRate}(p))$\;
    }
    $Selected \leftarrow \text{Top}(Peers, Score, target\_fanout)$\;
    \ForEach{$p \in Selected$}{
        Gửi $mesbt(b_i)$ tới $p$ qua TCP\;
    }
    CheckReconstruction()\;
}
\end{algorithm}

\textbf{Giai đoạn 3 (Leaderless Block Supplement):} Để loại bỏ điểm lỗi đơn duy nhất (Single Point of Failure - vai trò điều phối viên trung tâm), EdgeHydra áp dụng cơ chế bù đắp dữ liệu phi tập trung không cần leader. Mỗi trạm biên định kỳ 100ms sẽ phát quảng bá (broadcast) một gói tin Heartbeat UDP có kiểu bản tin là meshb tới toàn mạng với cấu trúc: \{"sender": "edge\_x", "received\_block\_ids": [0, 2]\} để thông báo danh sách các block ID mà nó đã thu thập thành công. Khi trạm biên e\_y nhận được heartbeat từ e\_x, nó đối chiếu với danh sách block hiện có của mình. Nếu phát hiện e\_x đang sở hữu block dữ liệu mà e\_y còn thiếu, e\_y sẽ chủ động thiết lập kết nối TCP và gửi yêu cầu kéo dữ liệu qua thông điệp mesbrq. Trạm biên e\_x nhận được yêu cầu sẽ phản hồi và gửi khối dữ liệu tương ứng cho e\_y qua thông điệp mesbs. Để chống nghẽn và lặp dữ liệu, hệ thống triển khai cơ chế Pacing (cooldown giữa các lần yêu cầu) và sổ cái gửi tin (\_stage2\_sent\_ledger, \_supplement\_sent\_ledger) để ngăn việc gửi trùng lặp cùng một khối dữ liệu trên một kênh truyền.

\begin{algorithm}[H]
\caption{Bù đắp dữ liệu phi tập trung (Stage 3)}
\textbf{Sự kiện Timer (mỗi 100ms):}\\
Broadcast $meshb(\text{owned\_blocks})$ tới toàn mạng\;
\vspace{2mm}
\textbf{Khi nhận được } $meshb(blocks\_list)$ \textbf{ từ } $Peer_x$: \\
\ForEach{$b_k \in blocks\_list$}{
    \If{$b_k \notin \text{owned\_blocks}$}{
        Gửi yêu cầu $mesbrq(b_k)$ tới $Peer_x$\;
        \textbf{break}\;
    }
}
\vspace{2mm}
\textbf{Khi nhận được } $mesbrq(b_k)$ \textbf{ từ } $Peer_y$: \\
\If{$b_k \in \text{owned\_blocks}$ và hết thời gian Cooldown}{
    Gửi cục bộ $mesbs(b_k)$ tới $Peer_y$\;
}
\end{algorithm}

\textbf{Giai đoạn 4 (Reconstruction \& Stop Signal):} Ngay khi một trạm biên tích lũy đủ k = 3 khối dữ liệu bất kỳ trong tổng số m = 4 khối (từ bất kỳ nguồn nào: trực tiếp từ Cloud ở Stage 1, chuyển tiếp ở Stage 2, hay bù đắp ở Stage 3), trạm biên đó sẽ gọi hàm giải mã của thư viện zfec để tái cấu trúc nguyên bản tệp tin gốc. Tệp tin khôi phục được lưu trữ vào thư mục kết quả. Đồng thời, trạm biên phát đi thông điệp dừng truyền mesdr\_bcast tới tất cả các trạm biên khác trong mạng lưới: \{"sender": "edge\_y", "reconstructed": true\}. Các trạm biên còn lại khi nhận được thông điệp này sẽ lập tức loại bỏ node e\_y ra khỏi danh sách hàng đợi gửi dữ liệu, giúp giải phóng băng thông mạng biên ngay lập tức.

\begin{algorithm}[H]
\caption{Giải mã và Tín hiệu dừng (Stage 4)}
\Fn{\text{CheckReconstruction}()}{
    \If{$|\text{owned\_blocks}| \geq k$}{
        $F_{recovered} \leftarrow \text{ErasureDecode}(\text{owned\_blocks}, k)$\;
        Lưu $F_{recovered}$ xuống ổ đĩa\;
        Broadcast $mesdr\_bcast$ tới tất cả các peer\;
        Gửi $mesdr$ về Cloud Server\;
    }
}
\vspace{2mm}
\textbf{Khi nhận được } $mesdr\_bcast$ \textbf{ từ } $Peer_z$: \\
Xóa $Peer_z$ khỏi mọi hàng đợi truyền tải\;
\end{algorithm}

\subsection{Thuật toán lập lịch truyền thích ứng}

\begin{figure}[H]
    \centering
    % Nhớ thay thế "images/adaptive.png" bằng đường dẫn thật tới ảnh bạn đã chụp
    \includegraphics[width=0.8\textwidth]{images/adaptive.png}
    \caption{Lưu đồ thuật toán lập lịch truyền thích ứng và kiểm soát độ phủ}
    \label{fig:adaptive}
\end{figure}

Để đo đạc chất lượng kết nối mạng biên theo thời gian thực, tiến trình giám sát mạng NetworkStateMonitor định kỳ gửi các gói tin thăm dò mesping qua TCP tới các peer lân cận và tính toán thời gian khứ hồi RTT thực tế. Nhằm loại bỏ các giá trị nhiễu đột biến, độ trễ RTT được làm mịn bằng thuật toán trung bình trượt lũy thừa (Exponential Weighted Moving Average - EWMA) theo công thức:

\[\mathrm{RTT}_{new}=\alpha\times\mathrm{RTT}_{measured}+(1-\alpha)\times\mathrm{RTT}_{old}\]

Trong đó, hệ số làm mịn được cấu hình là \ensuremath{\alpha} = 0.25. Điểm số đánh giá hiệu năng (Score) của mỗi peer được xác định dựa trên cả độ trễ RTT và tỷ lệ truyền tin thành công (Success Rate):

\[\mathrm{Score}=\mathrm{RTT}_{EWMA}+5000\times(1-\mathrm{Success\ Rate})\]

Các peer có kết nối kém hoặc bị sập sẽ có tỷ lệ thành công thấp hoặc RTT lớn, khiến điểm số tăng vọt. Trạm biên sẽ sắp xếp danh sách các peer theo điểm số tăng dần và luôn ưu tiên chuyển tiếp dữ liệu đến các node có điểm số thấp nhất.

\subsection{Mô hình chi phí truyền tải}

Chi phí truyền tải dữ liệu của hệ thống được tính toán và chuẩn hóa về đơn vị chi phí (Cost Units). Dựa trên thiết lập thực tế trong mã nguồn, chi phí băng thông truyền dữ liệu từ đám mây xuống trạm biên (Cloud-to-Edge - C2E) có giá đắt gấp 9 lần chi phí trao đổi mạng nội bộ giữa các trạm biên với nhau (Edge-to-Edge - E2E). Công thức tính toán tổng chi phí truyền tải cho một tệp tin là:

\[\mathrm{Cost}=\frac{\mathrm{Bytes}_{C2E}}{1{,}048{,}576}\times1.0+\frac{\mathrm{Bytes}_{E2E}}{1{,}048{,}576}\times\frac{1}{9}\]

Công thức này quy đổi toàn bộ dữ liệu truyền tải (tính bằng Byte) về Megabyte (MiB), trong đó 1 MiB truyền tải từ Cloud tương đương với 1 đơn vị chi phí, và 1 MiB truyền tải ngang hàng tại biên tương đương với 1/9 đơn vị chi phí.

\section{Đánh giá thực nghiệm}

\subsection{Thiết lập môi trường}

Hệ thống mô phỏng phân tán được triển khai và đánh giá trên môi trường Docker Desktop chạy trên hệ điều hành Windows 11. Mạng lưới giả lập bao gồm 5 container độc lập: 1 máy chủ đám mây (Cloud Node) và 4 máy chủ biên (Edge Nodes từ e\ensuremath{_0} đến e\ensuremath{_3}), liên kết với nhau thông qua mạng nội bộ Docker bridge network để mô phỏng mạng LAN tại biên.

Đối với giao thức EdgeHydra đề xuất, hệ thống áp dụng cấu hình mã hóa xóa k = 3, m = 4. Đối với giao thức baseline EdgeDis, tệp tin gốc được chia nhỏ thành n = 4 phân đoạn tĩnh bằng nhau, trong đó trạm biên e\ensuremath{_0} đóng vai trò là điều phối viên trung tâm (Coordinator). Cả hai giao thức đều được đo lường dưới cùng một điều kiện phần cứng và kích thước tệp tin kiểm thử mặc định là 24MiB.

\subsection{Các kịch bản thử nghiệm và mô phỏng lỗi (Fault Simulator)}

Để đánh giá một cách khách quan tính bền vững và khả năng chịu lỗi của EdgeHydra so với EdgeDis trong môi trường mạng biên biến động thực tế, nhóm tác giả đã tự thiết kế và lập trình tích hợp một Bộ giả lập sự cố (Fault Simulator) chuyên dụng. Bộ giả lập này hoạt động bằng cách tiêm lỗi (fault injection) trực tiếp vào tầng mạng TCP của các container thông qua các biến môi trường (Environment Variables) trong cấu hình Docker Compose (như \texttt{E2E\_FAULT\_DELAY\_MIN\_MS} hay \texttt{E2E\_FAULT\_CRASH\_AFTER\_S}). Thiết kế này cho phép giả lập chính xác 4 kịch bản thử nghiệm khắt khe nhất:

\textbf{Kịch bản 1 (Normal - Điều kiện mạng tiêu chuẩn):} Không kích hoạt bất kỳ cấu hình lỗi nào. Tất cả các trạm biên (từ e\ensuremath{_0} đến e\ensuremath{_3}) đều hoạt động với tốc độ truyền tải ngang hàng tối đa và tính toán ổn định. Kịch bản này đóng vai trò làm thước đo cơ sở (baseline) để so sánh chi phí và độ trễ ở điều kiện lý tưởng.

\textbf{Kịch bản 2 (Slow Servers - Suy thoái mạng cục bộ):} Kích hoạt độ trễ mạng ngẫu nhiên dao động liên tục từ 20ms đến 40ms trên 3 trạm biên (e\ensuremath{_1}, e\ensuremath{_2}, e\ensuremath{_3}) nhằm mô phỏng hiện tượng nghẽn mạng (Network Congestion) hoặc suy giảm chất lượng sóng không dây tại biên. Trạm biên e\ensuremath{_0} vẫn được giữ nguyên tốc độ lý tưởng. Kịch bản này thử thách độ nhạy bén của thuật toán lập lịch thích ứng (Adaptive Scheduling) trong việc phát hiện và né tránh các node bị chậm.

\textbf{Kịch bản 3 (Failed Servers - Sập nguồn đột ngột):} Giả lập lỗi nghiêm trọng nhất ở hạ tầng mạng biên: lỗi phần cứng hoặc mất điện đột ngột (Node Crash). Bộ giả lập ép buộc tiến trình máy chủ của trạm biên e\ensuremath{_1} phải bị ngắt kết nối hoàn toàn (kill process) đúng 1 giây kể từ khi hệ thống bắt đầu phân phối dữ liệu. Kịch bản này kiểm chứng khả năng tự bù đắp dữ liệu phi tập trung (Fault Tolerance) của mã hóa xóa mà không cần điểm điều phối trung tâm phát lệnh.

\textbf{Kịch bản 4 (Variable File Sizes - Đánh giá khả năng mở rộng):} Đánh giá năng lực xử lý (Scalability) của hệ thống khi phải truyền tải các tệp tin có kích thước phình to liên tục từ 1 MiB lên đến 144 MiB (thực hiện trong điều kiện mạng Normal). Mục tiêu là quan sát xem kiến trúc Leaderless của EdgeHydra có gặp phải hiện tượng nghẽn cổ chai (bottleneck) khi xử lý luồng dữ liệu cực lớn so với cấu trúc tập trung Coordinator của EdgeDis hay không.

\subsection{Kết quả thực nghiệm}

Dựa trên các kịch bản lỗi mạng đã thiết lập, chúng tôi tiến hành đo đạc và so sánh hiệu năng của giao thức EdgeHydra đề xuất với giao thức baseline EdgeDis. Bảng \ref{tab:fault_scenarios} trình bày tổng hợp kết quả đo lường về thời gian phân phối (Distribution Time - độ trễ hoàn thành tính bằng giây) và chi phí chuẩn hóa (Cost Units).

\begin{table}[H]
\centering
\caption{So sánh hiệu năng giữa EdgeHydra và EdgeDis (Kích thước tệp: 24MiB)}
\label{tab:fault_scenarios}
\small
\begin{tabularx}{\textwidth}{>{\raggedright\arraybackslash}X>{\raggedright\arraybackslash}X>{\raggedright\arraybackslash}X>{\raggedright\arraybackslash}X>{\raggedright\arraybackslash}X>{\raggedright\arraybackslash}X}
\toprule
\textbf{Kịch bản lỗi} & \textbf{Giao thức} & \textbf{Độ trễ (s)} & \textbf{C2E (MiB)} & \textbf{E2E (MiB)} & \textbf{Tổng chi phí (Units)} \\
\midrule
Normal (Mạng bình thường) & EdgeDis & 2.71 & 24.0 & 86.0 & 33.56 \\
 & \textbf{EdgeHydra} & \textbf{0.22} & 32.0 & 64.0 & \textbf{39.11} \\
\addlinespace
Slow Servers (Nghẽn mạng) & EdgeDis & 2.59 & 24.0 & 92.0 & 34.22 \\
 & \textbf{EdgeHydra} & \textbf{0.26} & 32.0 & 64.0 & \textbf{39.11} \\
\addlinespace
Failed Servers (Sập node đột ngột) & EdgeDis & 2.13 & 24.0 & 98.0 & 34.89 \\
 & \textbf{EdgeHydra} & \textbf{0.20} & 32.0 & 48.0 & \textbf{37.33} \\
\bottomrule
\end{tabularx}
\end{table}

Dữ liệu từ Bảng \ref{tab:fault_scenarios} cho thấy sự chênh lệch rõ rệt về mặt hiệu năng giữa hai phương pháp. Trong tất cả các kịch bản, EdgeHydra đều có \textbf{độ trễ (Distribution Time) thấp hơn EdgeDis khoảng 10 lần}. Cụ thể:

\begin{itemize}
    \item Trong kịch bản \textbf{Normal}, EdgeDis tốn 2.71 giây để hoàn thành phân phối, do phải dựa vào node điều phối viên (Coordinator) giám sát tuần tự, tạo ra nút thắt cổ chai về mặt lý thuyết. Trong khi đó, EdgeHydra với cơ chế phân phối song song Erasure Coding và thông báo chéo qua Heartbeat chỉ tốn 0.22 giây.
    \item Trong các kịch bản \textbf{Slow Servers} và \textbf{Failed Servers}, khi một số node mạng bị nghẽn hoặc sập nguồn, độ trễ của EdgeDis không thay đổi nhiều (dao động từ 2.13s đến 2.59s), nhưng nguyên nhân là do lượng dữ liệu truyền tải Edge-to-Edge (E2E) tăng vọt lên tới 98.0 MiB để bù đắp các khối lỗi. Hệ thống EdgeDis bị chững lại để chờ điều phối viên phát hiện lỗi và yêu cầu truyền lại dữ liệu.
    \item Ngược lại, \textbf{EdgeHydra} cho thấy khả năng dung lỗi (Fault Tolerance) xuất sắc. Do áp dụng mô hình mã hóa xóa $k=3, m=4$, mỗi node chỉ cần nhận đủ 3 khối dữ liệu để khôi phục tệp 24MiB, hoàn toàn không cần chờ node thứ 4 bị sập. Việc áp dụng thuật toán lập lịch thích ứng (Adaptive Scheduling) giúp EdgeHydra chủ động điều hướng dữ liệu qua các đường truyền có chất lượng RTT tốt hơn, giảm thiểu lượng dữ liệu dư thừa trên mạng E2E xuống chỉ còn 48.0 MiB trong kịch bản lỗi node.
\end{itemize}

Tuy nhiên, sự tối ưu về tốc độ của EdgeHydra đi kèm với một sự gia tăng nhẹ về mặt chi phí băng thông (Cost). Lưu lượng C2E của EdgeHydra luôn cố định ở mức 32.0 MiB (cao hơn so với 24.0 MiB của EdgeDis) do tính chất tạo dữ liệu dư thừa từ quá trình mã hóa xóa ban đầu. Hệ quả là tổng chi phí (Total Cost Units) của EdgeHydra cao hơn EdgeDis từ 7\% đến 16\%. Mặc dù vậy, sự đánh đổi (trade-off) này là hoàn toàn chấp nhận được và mang tính thực tiễn cao đối với các dịch vụ mạng biên ưu tiên độ trễ siêu thấp (ultra-low latency). Bảng \ref{tab:e2e_c2e_overhead} minh họa cụ thể tỷ lệ gia tăng chi phí băng thông mạng nội bộ (E2E) so với mạng Cloud (C2E) của hai giao thức.

\begin{table}[H]
\centering
\caption{Tỷ lệ chi phí truyền tải C2E và E2E (Cost Overhead)}
\label{tab:e2e_c2e_overhead}
\small
\begin{tabularx}{0.8\textwidth}{>{\raggedright\arraybackslash}X>{\raggedright\arraybackslash}X}
\toprule
\textbf{Giao thức} & \textbf{Tỷ lệ E2E / C2E} \\
\midrule
\textbf{EdgeHydra} & Cố định $\times 2.00$ (ở mọi kích thước tệp) \\
EdgeDis & Biến thiên từ $\times 3.08$ đến $\times 4.00$ \\
\bottomrule
\end{tabularx}
\end{table}

\subsection{Đánh giá khả năng mở rộng (Scalability) với kích thước tệp biến đổi}

Để chứng minh khả năng mở rộng của thuật toán, chúng tôi tiến hành chạy \textbf{Kịch bản 4 (Variable File Sizes)} với kích thước tệp dữ liệu phân phối thay đổi từ 1 MiB đến 144 MiB trong điều kiện mạng bình thường. 

\begin{table}[H]
\centering
\caption{Thời gian phân phối theo kích thước tệp (Giây)}
\label{tab:variable_sizes}
\small
\begin{tabularx}{\textwidth}{>{\centering\arraybackslash}X>{\centering\arraybackslash}X>{\centering\arraybackslash}X>{\centering\arraybackslash}X}
\toprule
\textbf{Kích thước File} & \textbf{EdgeDis (s)} & \textbf{EdgeHydra (s)} & \textbf{Speedup} \\
\midrule
\textbf{1 MiB} & 1.53 & \textbf{0.12} & 12.7 $\times$ \\
\textbf{24 MiB} & 2.49 & \textbf{0.24} & 10.4 $\times$ \\
\textbf{48 MiB} & 3.78 & \textbf{0.36} & 10.5 $\times$ \\
\textbf{72 MiB} & 2.42 & \textbf{0.46} & 5.2 $\times$ \\
\textbf{96 MiB} & 2.78 & \textbf{0.57} & 4.9 $\times$ \\
\textbf{120 MiB} & 3.35 & \textbf{0.98} & 3.4 $\times$ \\
\textbf{144 MiB} & 4.11 & \textbf{0.88} & 4.7 $\times$ \\
\bottomrule
\end{tabularx}
\end{table}

Kết quả thu được từ Bảng \ref{tab:variable_sizes} cho thấy độ trễ của EdgeHydra tuân theo mô hình tăng trưởng tuyến tính và rất ổn định ($O(1)$ scaling behavior) khi kích thước tệp phình to. Ngay cả khi đối mặt với tải trọng lớn nhất là tệp 144 MiB, EdgeHydra cũng chỉ tiêu tốn 0.88 giây để hoàn thành phân phối toàn mạng. Ở chiều ngược lại, thuật toán baseline EdgeDis bộc lộ rõ sự thiếu ổn định. Thời gian phân phối của EdgeDis biến thiên khó đoán và đạt đỉnh điểm hơn 4.1 giây cho tệp 144 MiB. Lý do là khi kích thước phân đoạn lớn hơn, thời gian xử lý và kiểm tra tính toàn vẹn của mỗi khối dữ liệu cũng tăng theo, tạo áp lực lớn lên node Coordinator trong việc thu thập tín hiệu hoàn thành từ các node khác. 

Qua kịch bản này, chúng ta có thể khẳng định rằng kiến trúc chia sẻ ngang hàng không điều phối viên (Leaderless) kết hợp lập lịch thích ứng của EdgeHydra vượt trội hoàn toàn về mặt khả năng mở rộng với dữ liệu dung lượng lớn, đáp ứng tiêu chuẩn khắt khe của các hệ thống Edge Caching trong thực tế như phân phối mô hình AI hay luồng video băng thông cao. Nhằm cung cấp một góc nhìn tổng quan nhất về sự cân bằng giữa hiệu suất và chi phí, Bảng \ref{tab:tradeoff_score} tổng hợp Điểm số đánh đổi (Trade-off Score = Speedup / Cost Overhead) qua các kích thước tệp khác nhau.

\begin{table}[H]
\centering
\caption{Điểm số đánh đổi Tốc độ - Chi phí (Trade-off Score) theo kích thước tệp}
\label{tab:tradeoff_score}
\small
\begin{tabularx}{\textwidth}{>{\centering\arraybackslash}X>{\centering\arraybackslash}X>{\centering\arraybackslash}X>{\centering\arraybackslash}X}
\toprule
\textbf{Kích thước File} & \textbf{Tăng tốc (Speedup)} & \textbf{Phụ phí (Cost Overhead)} & \textbf{Điểm số (Trade-off)} \\
\midrule
1 MiB & $12.4\times$ & $1.15\times$ & 10.8 \\
24 MiB & $10.4\times$ & $1.13\times$ & 9.3 \\
48 MiB & $10.4\times$ & $1.21\times$ & 8.6 \\
72 MiB & $5.2\times$ & $1.21\times$ & 4.3 \\
96 MiB & $4.9\times$ & $1.20\times$ & 4.1 \\
120 MiB & $3.4\times$ & $1.20\times$ & 2.8 \\
144 MiB & $4.7\times$ & $1.20\times$ & 3.9 \\
\bottomrule
\end{tabularx}
\end{table}

\section{Hạn chế và Hướng phát triển}

Mặc dù giao thức EdgeHydra đã chứng minh được tính ưu việt vượt trội về độ trễ và khả năng chịu lỗi so với mô hình truyền thống, quá trình thực nghiệm và thiết kế hệ thống vẫn còn tồn tại một số hạn chế nhất định mở ra các hướng nghiên cứu trong tương lai:

\textbf{1. Chi phí mã hóa và giải mã (Encoding/Decoding Overhead):}
Giải pháp hiện tại sử dụng mã hóa xóa (Erasure Coding) giúp giảm băng thông mạng và loại bỏ điểm lỗi đơn. Tuy nhiên, quá trình tính toán phân rã file tại Cloud và đặc biệt là giải mã tái cấu trúc tệp tại trạm biên tốn khá nhiều tài nguyên CPU. Đối với các Edge Node có phần cứng giới hạn (như thiết bị IoT), độ trễ tính toán này có thể trở thành nút thắt cổ chai mới.
\textit{Hướng phát triển:} Nghiên cứu tích hợp các giải thuật mã hóa tối ưu mức phần cứng (như tận dụng tập lệnh SIMD hoặc tính toán đa luồng song song trên GPU) để giảm thiểu triệt để độ trễ xử lý.

\textbf{2. Chi phí mạng nội bộ (E2E Bandwidth Cost):}
Nhờ cơ chế chia sẻ bù đắp dữ liệu phi tập trung, EdgeHydra duy trì được tốc độ cao nhưng lại làm tăng chi phí băng thông E2E lên cố định khoảng $2.0\times$ so với mạng Cloud. Trong một số hạ tầng mạng có tính phí dung lượng nội bộ cao, sự đánh đổi này có thể là một rào cản.
\textit{Hướng phát triển:} Cải tiến thuật toán Lập lịch thích ứng (Adaptive Scheduling) bằng cách bổ sung thêm trọng số "chi phí đường truyền" (Link Cost) vào hàm tính Điểm (Score). Hệ thống sẽ tự động cân bằng giữa việc chọn đường truyền nhanh nhất và đường truyền rẻ nhất.

\textbf{3. Cấu hình tham số $(k, m)$ tĩnh:}
Đồ án hiện tại đang tiến hành đo đạc dựa trên cấu hình mã hóa cố định $k=3, m=4$. Trong thực tế, chất lượng mạng không dây tại biên luôn biến động theo thời gian. Việc giữ cố định cấu hình này có thể dẫn đến lãng phí tài nguyên mạng khi tín hiệu đang tốt, hoặc không đủ lượng dữ liệu bù đắp khi mạng suy thoái nghiêm trọng.
\textit{Hướng phát triển:} Thiết kế một mô hình Học tăng cường (Reinforcement Learning) nhằm giám sát trạng thái mạng tổng thể và tự động tinh chỉnh linh hoạt tỷ lệ $(k, m)$ theo thời gian thực (Dynamic Erasure Coding), qua đó luôn duy trì Điểm số đánh đổi (Trade-off Score) ở mức tối ưu nhất.

\section{Kết luận}

Trong nghiên cứu này, chúng tôi đã triển khai, đánh giá và tối ưu hóa giải pháp phân phối dữ liệu mạng biên EdgeHydra dựa trên nền tảng kỹ thuật mã hóa xóa (Erasure Coding). Bằng việc khắc phục các hạn chế chí mạng của mô hình truyền dữ liệu song song truyền thống (EdgeDis), EdgeHydra giúp hệ thống đạt được tốc độ truyền tải cực cao đồng thời miễn nhiễm với các lỗi sập nguồn đột ngột hay nghẽn mạng cục bộ.

Đặc biệt, đóng góp chính của nhóm nghiên cứu là việc đề xuất tích hợp \textbf{Thuật toán lập lịch truyền thích ứng (Adaptive Scheduling)} dựa trên độ trễ RTT và kiểm soát độ phủ mạng. Cải tiến này đã giải quyết triệt để vấn đề "Bão tin nhắn" (Broadcast Storm) trong cơ chế chia sẻ ngang hàng Leaderless nguyên bản, giúp giảm tải lượng dữ liệu dư thừa chạy trên mạng lưới giữa các trạm biên. Kết quả thực nghiệm thông qua các kịch bản giả lập trên nền tảng Docker đã chứng minh EdgeHydra mang lại tốc độ phân phối \textbf{nhanh hơn trung bình 10 lần} so với EdgeDis, chỉ đánh đổi một lượng chi phí mạng nội bộ không đáng kể (7\%). Trong tương lai, mô hình này có thể được mở rộng bằng cách tự động tinh chỉnh linh hoạt tỷ lệ mã hóa xóa $(k, m)$ theo thời gian thực dựa trên trạng thái học máy của hệ thống mạng.

\begin{thebibliography}{9}

\bibitem{edgedis} 
B. Li, Q. He, F. Chen, L. Lyu, A. Bouguettaya, and Y. Yang, "EdgeDis: Enabling fast, economical, and reliable data dissemination for mobile edge computing," \textit{IEEE Transactions on Services Computing}, vol. 17, no. 4, pp. 1504–1518, Jul./Aug. 2024.

\bibitem{erasurecoding} 
Q. He \textit{et al.}, "EdgeHydra: Fault-Tolerant Edge Data Distribution Based on Erasure Coding," \textit{IEEE Transactions on Parallel and Distributed Systems}, vol. 36, no. 1, pp. 29–42, Jan. 2025.

\end{thebibliography}
\end{document}

